\subsection{A simultaneous count of separated significant directions}
\label{count:section}

On the set $F_{\mathrm{conc}}$ from \Cref{conc:mask-removal}, the weighted
threshold \eqref{eq:src-12} bounds the total top mass by $NL^{C_n}$.
Because each significant top contributes at least one, this bounds their
number by the same quantity, which may grow with $N$.  For the later finite
Fourier gain we need a bound polynomial in $L$ for significant directions
that are $\rho$-separated inside an interval of length $B\rho$, with top
precisions at most $\rho/B^2$.  We derive that bound from the popularity
witnesses, simultaneously for every $\rho>0$ and every such list.

We retain the parameters and the finite family of tops from the preceding
sections.  In particular, every top has width $W=mw$, length
$l_T\leq\tau$, and precision $d_T=w/l_T$.  Its witness satisfies
\begin{equation}\label{count:witness-data}
 E_T\subset L^{a_0}R_T,\qquad
 |E_T|\geq L^{-a_0}l_TW,\qquad
 |u-s_T|\leq c_0d_T\quad\hbox{on }E_T,
\end{equation}
where $a_0$ and $c_0$ are fixed constants; increasing $a_0$ absorbs
the fixed constants in the first two bounds.  The dilated box in
\eqref{count:witness-data} is understood in the coordinates of $R_T$.
The sum of the top areas is $S=\sum_T l_TW$.  A top is significant at
$z$ when $\dist_T(z)\leq K_0$, with $K_0=L^{A+2}$ as before.

\begin{proposition}\label{count:simultaneous}
There is a fixed exponent $C_D$ with the following property.  For every
fixed $J>0$ and all sufficiently large $L$, there is a measurable set
$\mathcal E_{\mathrm{count}}\subset\R^2$ such that
\begin{equation}\label{count:exceptional-budget}
 |\mathcal E_{\mathrm{count}}|\leq S e^{-JL}.
\end{equation}
Define
\begin{equation}\label{count:good-set}
 Z_0=F_{\mathrm{conc}}\setminus\mathcal E_{\mathrm{count}},
 \qquad D=L^{C_D}.
\end{equation}
For every $z\in Z_0$, every $\rho>0$, and every list $\mathcal L$ of
tops significant at $z$ satisfying
\begin{equation}\label{count:list-hypotheses}
 \begin{split}
 &\max_{T\in\mathcal L}s_T-\min_{T\in\mathcal L}s_T\leq B\rho,
 \qquad d_T\leq\rho/B^2\quad(T\in\mathcal L),\\
 &|s_T-s_{T'}|\geq\rho\quad(T,T'\in\mathcal L,\ T\ne T'),
 \end{split}
\end{equation}
one has $\#\mathcal L\leq D/100$.
The exponent $C_D$ depends only on the fixed geometric constants,
including the already chosen mask exponent $A$.  It is independent of
the number of tops, the number of scales, and the bin Hilbert space.
\end{proposition}

The margin in the conclusion permits a fixed number of additional
directions and fixed multiplicity factors in later applications, after
increasing the absolute lower bound on $L$.  The exceptional set will
be constructed from one angular bin family for each reference top.
It therefore applies to all $\rho$ and all lists simultaneously.

We transfer the other tops' popularity witnesses to a line in the direction
of a longest top.  Separated top directions then occupy distinct angular
bins, and the following lemma bounds the number of bins each having a
large average on some interval through the evaluation point.
The reference top contributes one additional direction.

\subsection{A counting lemma on a line}

The popular-bin counting argument is inspired by Bateman's one-variable
covering/BMO method \cite[Section~3.3.2, Lemmas~5--6]{Bateman2008Maximal}.
The moment estimate below is proved directly; the Lipschitz transfer to
reference lines and the simultaneous multiresolution formulation are
established here, rather than imported from Bateman's maximal theorem.

\begin{lemma}\label{count:line-lemma}
Let $(f_b)_{b\in\mathcal B}$ be a countable family of measurable
functions on $\R$ such that $0\leq f_b\leq1$ and
$\sum_b f_b\leq C_0$, where $C_0\geq1$.  Fix $H>0$ and
$0<\vartheta\leq1$.  There is a measurable function
$\mathcal N:\R\to[0,\infty]$ with the following properties.
If $\mathcal B_x$ is any collection of distinct indices such that,
for every $b\in\mathcal B_x$, some interval $I_b$ containing $x$
satisfies
\begin{equation}\label{count:line-average}
 0<|I_b|\leq H,\qquad |I_b|^{-1}\int_{I_b}f_b\geq\vartheta,
\end{equation}
then $\#\mathcal B_x\leq\mathcal N(x)$.  Moreover, for every bounded
interval $U$ and every $\lambda\geq0$,
\begin{equation}\label{count:line-tail}
 |\{x\in U:\mathcal N(x)>\lambda\}|
 \leq C(|U|+H)\exp(-c\vartheta\lambda/C_0),
\end{equation}
where $C,c>0$ are absolute.  If the $f_b$ depend jointly measurably
on an additional parameter, $\mathcal N$ can be chosen jointly
measurable in that parameter and $x$.
\end{lemma}

\begin{proof}
For $t\in\{0,1/3,2/3\}$ use the dyadic grid whose intervals of
length $2^k$ are
\[
 \mathcal D_k^t
 =\big\{2^k[j+(-1)^kt,j+1+(-1)^kt):j\in\Z\big\}.
\]
These grids are nested: a boundary at scale $k+1$ is a boundary
at scale $k$ because $3t$ is an integer.  At a fixed scale their
combined boundary points have spacing $2^k/3$.  Given an interval
$I$ of positive finite length, choose the least integer $k$ with
$2^k>3|I|$.  Its closure meets boundary points from at most one of
the three grids.  Hence $I$ is contained in an interval $Q$ from
one of the other grids, and $|Q|\leq6|I|$.

Choose an integer $k_H$ with $6H\leq2^{k_H}<12H$, and in each grid
retain only scales $k\leq k_H$.  Put $\beta=\vartheta/6$.  For each
$b$ and each grid, call $Q$ bad if
\[
 |Q|^{-1}\int_Q f_b\geq\beta.
\]
Let $\mathcal J_b^t$ be the intervals maximal under inclusion among
these bad intervals.  Every bad interval is contained in a maximal
one, since it has only finitely many ancestors up to level $k_H$.
For fixed $b,t$, the members of $\mathcal J_b^t$ are pairwise
disjoint.  Define
\[
 \mathcal N_t(x)=\sum_{b\in\mathcal B}
                   \sum_{Q\in\mathcal J_b^t}\one_Q(x),
 \qquad \mathcal N(x)=\sum_{t\in\{0,1/3,2/3\}}\mathcal N_t(x).
\]
Each interval in \eqref{count:line-average} produces a bad dyadic
interval containing $x$, at an allowed scale.  Thus each distinct
index in $\mathcal B_x$ contributes at least once to $\mathcal N(x)$.

For every dyadic interval $P$ at an allowed scale, disjointness
within each bin gives the Carleson estimate
\begin{equation}\label{count:carleson}
 \begin{split}
 \sum_b\sum_{\substack{Q\in\mathcal J_b^t\\Q\subseteq P}}|Q|
 &\leq\beta^{-1}\sum_b
            \sum_{\substack{Q\in\mathcal J_b^t\\Q\subseteq P}}
                         \int_Q f_b\\
 &\leq\beta^{-1}\int_P\sum_b f_b
 \leq\mathcal A|P|,\qquad \mathcal A=C_0/\beta.
 \end{split}
\end{equation}
All sums here are of nonnegative terms, so this argument also
applies to countably many bins and intervals.

For completeness, \eqref{count:carleson} has an elementary
exponential-tail consequence.  Fix a root interval $P_0$ of length
$2^{k_H}$.  First take any finite subfamily of the selected interval
occurrences inside $P_0$, retaining multiplicities when different
bins select the same interval, and let $V$ count that subfamily.
For an integer $r\geq1$, expansion of $V^r$ and the nesting of
intersecting dyadic intervals give
\begin{equation}\label{count:moment}
 \int_{P_0} V^r
 \leq r!\sum_{Q_1\supseteq\cdots\supseteq Q_r}|Q_r|
 \leq r!\mathcal A^r|P_0|.
\end{equation}
The chain sum counts occurrences with their multiplicities.
To obtain the second inequality, sum first over $Q_r\subseteq
Q_{r-1}$ and apply \eqref{count:carleson}; repeat until the
remaining sum is bounded by $\mathcal A|P_0|$.  Repeated intervals
are allowed throughout.  Ordering the $r$ occurrences accounts
for the factor $r!$ and can only overcount.

Enumerate all selected occurrences inside $P_0$.  Their finite
initial sums increase to $\mathcal N_t$ on $P_0$, so monotone
convergence and \eqref{count:moment} show
\begin{equation}\label{count:exponential-moment}
 \int_{P_0}\exp(\mathcal N_t/(2\mathcal A))
 \leq \sum_{r=0}^{\infty}2^{-r}|P_0|=2|P_0|.
\end{equation}
In particular the count is finite almost everywhere.  Markov's
inequality bounds its superlevel measure in $P_0$ by
$2|P_0|\exp(-\lambda/(2\mathcal A))$.
The root intervals of one grid meeting $U$ have total length at
most $|U|+2^{k_H+1}$.  Sum the root estimates and observe that
$\mathcal N>\lambda$ implies $\mathcal N_t>\lambda/3$ for at least
one grid.  Since $\mathcal A=6C_0/\vartheta$ and $2^{k_H}<12H$,
this proves \eqref{count:line-tail}.

There are countably many bins and dyadic intervals.  For each
fixed interval its badness is determined by a measurable integral.
Its maximality adds only finitely many conditions on its ancestors
before level $k_H$.  Consequently the displayed count is
measurable.  The same argument, using measurability of integrals
with a parameter, proves the final assertion.  No maximality
selection or exceptional set depending on an uncountable family
of intervals is required.
\end{proof}

\subsection{Transferring witnesses to a reference line}

Fix a top $R$.  We use the coordinates
\[
 z=(x,y_R+s_R(x-x_R)+v).
\]
The significant region of $R$ is contained in
\begin{equation}\label{count:reference-box}
 \mathcal Q_R=\{|x-x_R|\leq K_0l_R,\ |v|\leq K_0W\}.
\end{equation}
The coordinate change $(x,v)\mapsto z$ has Jacobian one.

For each integer $h$, partition each of the two angular annuli
\[
 [s_R+2^h,s_R+2^{h+1}),\qquad
 (s_R-2^{h+1},s_R-2^h]
\]
into $B^2$ intervals of length $\delta_h=2^h/B^2$.
Here $B$ is a dyadic integer, as stipulated in its definition.
Denote this countable family of base intervals by
$\mathcal B_R$.  Enlarge every base interval concentrically by
a fixed factor $c_1$, chosen below, and write $b^+$ for the
resulting closed interval.  For sufficiently large $L$, and hence
$B$, these enlarged intervals satisfy
\begin{equation}\label{count:bin-overlap}
 \sum_{b\in\mathcal B_R}\one_{b^+}(a)\leq C_{\mathrm{bin}}
 \qquad(a\in\R),
\end{equation}
with a fixed $C_{\mathrm{bin}}$.  Indeed, enlargement moves an
endpoint by at most $c_1 2^h/B^2$.  Once $B^2\geq4c_1$, an
enlarged interval from annulus $h$ stays at distance comparable
to $2^h$ from $s_R$.  A point therefore meets only a fixed number
of annuli, and at most a fixed number, depending on $c_1$, of
the enlarged equal-length intervals in each such annulus.

Choose a fixed power $P_L=L^{a_1}$ sufficiently large to dominate
the dilation constants in \eqref{count:witness-data}, $K_0$,
and their fixed multiples.  It will not depend on $R$ or on a
list.  Replacing a witness for this argument by a compact subset
$\widetilde E_T\subset E_T$ of at least half its measure is
permitted by inner regularity.  The slope condition can first
be enforced after deleting a null subset if necessary.  The
horizontal projection $A_T$ of $\widetilde E_T$ is compact, and
Fubini's theorem gives
\begin{equation}\label{count:projection-size}
 |A_T|\geq \frac{|\widetilde E_T|}{2L^{a_0}W}
          \geq \frac14L^{-2a_0}l_T.
\end{equation}
The denominator bounds the length of each vertical section of
the dilated witness box.  The original witnesses themselves
remain unchanged elsewhere in the proof.

Suppose now that a list satisfies \eqref{count:list-hypotheses}
at $z$, and choose $R$ from the list with maximal length.  If
the list is empty there is nothing to prove.  For $T\ne R$ put
$\Delta_T=|s_T-s_R|$.  Then
\begin{equation}\label{count:relative-precision}
 \rho\leq\Delta_T\leq B\rho,\qquad
 d_T\leq\Delta_T/B^2,\qquad l_T\leq l_R\leq\tau.
\end{equation}
There is an interval $I_T$ containing both $A_T$ and the horizontal
coordinate $x$ of $z$ with
\begin{equation}\label{count:witness-interval}
 |I_T|\leq2P_Ll_T\leq2P_Ll_R.
\end{equation}
For example use the interval centered at $x_T$ of radius
$P_Ll_T$; significance gives $|x-x_T|\leq K_0l_T$.

For $x'\in A_T$, take any $(x',y')\in\widetilde E_T$, and let
$z'=(x',y+s_R(x'-x))$ be the point with horizontal coordinate
$x'$ on the reference line through $z$.  The top-coordinate
representations of the witness point and $z$ give
\begin{equation}\label{count:vertical-distance}
 |y'-(y+s_R(x'-x))|
 \leq P_L\bigl(W+l_T\Delta_T\bigr).
\end{equation}
This bound holds for every possible witness point above $x'$,
so no measurable choice of such points is needed.  By the
Lipschitz bound on $u$,
\begin{equation}\label{count:transferred-error}
 |u(z')-s_T|
 \leq c_0d_T+\Lip(u)P_L\bigl(W+l_T\Delta_T\bigr).
\end{equation}
The parameter choices ensure
\begin{equation}\label{count:small-time}
 \Lip(u)P_Lm\tau\leq B^{-100},\qquad B^2\leq m
\end{equation}
for sufficiently large $L$.  To check the first assertion,
$\log(P_Lm\tau)= -L^{0.97}+O(L^{0.9}+\log L)$, whereas
$\log B=O(L^{0.7}(\log L)^2)$; the fixed Lipschitz constant
does not affect this comparison.  Since
$W=ml_Td_T\leq m\tau\Delta_T/B^2$, the last two terms in
\eqref{count:transferred-error} are at most
\[
 \bigl(\Lip(u)P_Lm\tau+
       \Lip(u)P_L\tau B^2\bigr)\frac{\Delta_T}{B^2}
 \leq 2B^{-100}\frac{\Delta_T}{B^2}.
\]
Consequently the transferred error is bounded by
$(c_0+1)\Delta_T/B^2$.

Let $b(T)$ be the unique base interval containing $s_T$.  If its
annulus index is $h$, then $2^h\leq\Delta_T<2^{h+1}$.
Choose $c_1$ large enough in terms of $c_0$ so that every point
within $2(c_0+1)\delta_h$ of any point of a base interval lies
in its enlargement.  Equations~\eqref{count:transferred-error}
and \eqref{count:small-time} now show
\[
 u(z')\in b(T)^+\qquad(x'\in A_T).
\]
Define, on every reference line, the bin functions
\begin{equation}\label{count:line-bin-functions}
 f_{R,b}(x',v)=
 \one_{b^+}\bigl(u(x',y_R+s_R(x'-x_R)+v)\bigr).
\end{equation}
They are jointly measurable, and their sum is bounded by
$C_{\mathrm{bin}}$ by \eqref{count:bin-overlap}.
Combining \eqref{count:projection-size} and
\eqref{count:witness-interval}, there is a fixed exponent $a_2$
such that, with $\vartheta=L^{-a_2}$,
\begin{equation}\label{count:popular-line-bin}
 \frac1{|I_T|}\int_{I_T} f_{R,b(T)}(x',v)\,\dd x'
 \geq\vartheta.
\end{equation}
For example any fixed $a_2>2a_0+a_1$ suffices once $L$ is large.
All choices so far precede the eventual choice of $D$.

The base intervals $b(T)$ for $T\in\mathcal L\setminus\{R\}$
are distinct.  Indeed, any one such interval has length
$2^h/B^2\leq\Delta_T/B^2\leq\rho/B<\rho$, so it cannot
contain two directions from a $\rho$-separated list.

\subsection{The common exceptional set and the choice of \texorpdfstring{$D$}{D}}

Apply \Cref{count:line-lemma} to the full family
\eqref{count:line-bin-functions}, with
\[
 H_R=2P_Ll_R,\qquad \vartheta=L^{-a_2},\qquad
 C_0=C_{\mathrm{bin}}.
\]
Write $\mathcal N_R(x,v)$ for its jointly measurable count.
This construction depends on $R$ but on neither $\rho$ nor a
list.  Whenever $R$ is a longest member of an admissible list
at $z=(x,y_R+s_R(x-x_R)+v)$, the preceding argument gives
\begin{equation}\label{count:list-by-one-count}
 \#\mathcal L\leq1+\mathcal N_R(x,v).
\end{equation}
Use the threshold $\Lambda_{\mathrm{count}}=L^{a_2+2}$ and let
\[
 \mathcal E_R=
 \{z\in\mathcal Q_R:
              \mathcal N_R(x,v)>\Lambda_{\mathrm{count}}\}.
\]
In \eqref{count:line-tail}, take
$U=[x_R-K_0l_R,x_R+K_0l_R]$.  Since $P_L$ dominates $K_0$,
that estimate gives, for each $v$,
\[
 |\{x\in U:\mathcal N_R(x,v)>
                         \Lambda_{\mathrm{count}}\}|
 \leq CP_Ll_R e^{-cL^2}.
\]
Integrating over $|v|\leq K_0W$ and using the Jacobian-one
coordinates yields
\begin{equation}\label{count:one-top-exception}
 |\mathcal E_R|\leq CP_LK_0l_RW e^{-cL^2}.
\end{equation}
In particular, for every fixed $J>0$, the right-hand side is
at most $e^{-JL}l_RW$ for sufficiently
large $L$, because $P_LK_0$ is a fixed power of $L$.

Set $\mathcal E_{\mathrm{count}}=\bigcup_R\mathcal E_R$.
The top family is finite, and the sets are measurable.  Summing
\eqref{count:one-top-exception} proves
\eqref{count:exceptional-budget}.  The same proof applies to
any countable family with finite area sum by countable
subadditivity.  The angular families were already countable,
and their finite approximations were handled in the proof of
\Cref{count:line-lemma}; thus no cardinality or depth factor is
introduced by a limiting operation.

Choose
\begin{equation}\label{count:choice-D}
 C_D=a_2+4,\qquad D=L^{C_D}.
\end{equation}
At a point outside $\mathcal E_{\mathrm{count}}$, a longest
member $R$ of any nonempty admissible list is significant, so
the point lies in $\mathcal Q_R$ and
$\mathcal N_R\leq\Lambda_{\mathrm{count}}$ there.
Equation~\eqref{count:list-by-one-count} gives
$\#\mathcal L\leq1+L^{a_2+2}\leq D/100$ for sufficiently
large $L$.  This proves \Cref{count:simultaneous}.

Finally, this additional discard is permissible in the original
adjoint pairing.  The ordinary model estimate gives
$\norm{P}_2\leq L^C\sqrt{|E|}$, and $S\leq L^C|E|$.
Consequently
\begin{equation}\label{count:discard-cost}
 \int_{F_{\mathrm{conc}}\setminus Z_0}\norm{P}
 \leq L^C|E|e^{-JL/2}
 \leq L^C e^{(2-J/2)L}\,N|F|,
\end{equation}
where the last inequality uses $|E|=N^2|F|$ and $N\leq e^{2L}$.
Taking, for example, any fixed $J>4$ makes this smaller than
any prescribed inverse power of $L$ times $N|F|$, after
increasing the absolute lower bound on $L$.  The estimates
\eqref{eq:src-11} and \eqref{eq:src-12}, together with the
already established mask replacement, continue to hold on
$Z_0\subset F_{\mathrm{conc}}$.
